% GENERATED FILE. Edit canonical lessons, then run npm run generate:books.
% canonical-source-hash: fnv1a64:57cfa4013da75e3e
% canonical-lessons: GU-C132-cokhkhu, GU-C132-gandu, GU-C132-khatu, GU-C132-tikhu, GU-C132-naram

\chapter{Clean, Dirty, Sour, Spicy, Soft}
\label{ch:gu-clean-dirty-sour-spicy-soft}
\hlchaptermodality{132}

\begin{chapteropening}
\textbf{By the end of this chapter:} \emph{I can say clean, dirty, sour, spicy and soft.}

Complete the last lesson of chapter 132: I can say clean, dirty, sour, spicy and soft.
\end{chapteropening}

\section[\texorpdfstring{\gu{ચોખ્ખું} (cokhkhũ)}{ચોખ્ખું (cokhkhũ)}]{\gu{ચોખ્ખું} (cokhkhũ) — clean}

\label{lesson:GU-C132-cokhkhu}



\begin{quote}
Before the new one: say the Gujarati for wet, then the Gujarati for dry.
\end{quote}

\subsection*{You'll want to know: \gu{ચોખ્ખું}}
\textbf{\gu{ચોખ્ખું}} (\emph{cokhkhũ}) — \textquotedblleft{}clean\textquotedblright{}.

\begin{grammarlens}[title={What you've built}]
One more word for this chapter.
\end{grammarlens}

\subsection*{Your turn}
\begin{itemize}
\raggedright
  \item \emph{Say it:} \emph{cokhkhũ}
  \item \emph{Say it:} \emph{cokhkhũ}, once more
  \item \emph{Say it:} say \emph{cokhkhũ}
  \item \emph{Recall:} say the Gujarati for wet, then the Gujarati for dry, then say \emph{cokhkhũ} again
\end{itemize}

\begin{tcolorbox}[breakable,colback=teal!4,colframe=teal!35!black,title={Before you move on}]
What does \gu{ચોખ્ખું} mean? (\textquotedblleft{}Clean\textquotedblright{}.) Say it once more.
\end{tcolorbox}
\section[\texorpdfstring{\gu{ગંદું} (gandũ)}{ગંદું (gandũ)}]{\gu{ગંદું} (gandũ) — dirty}

\label{lesson:GU-C132-gandu}



\begin{quote}
Before the new one: say the Gujarati for dry, then the Gujarati for clean.
\end{quote}

\subsection*{You'll want to know: \gu{ગંદું}}
\textbf{\gu{ગંદું}} (\emph{gandũ}) — \textquotedblleft{}dirty\textquotedblright{}.

\begin{grammarlens}[title={What you've built}]
One more word for this chapter.
\end{grammarlens}

\subsection*{Your turn}
\begin{itemize}
\raggedright
  \item \emph{Say it:} \emph{gandũ}
  \item \emph{Say it:} \emph{gandũ}, once more
  \item \emph{Say it:} say \emph{gandũ}
  \item \emph{Recall:} say the Gujarati for dry, then the Gujarati for clean, then say \emph{gandũ} again
\end{itemize}

\begin{tcolorbox}[breakable,colback=teal!4,colframe=teal!35!black,title={Before you move on}]
What does \gu{ગંદું} mean? (\textquotedblleft{}Dirty\textquotedblright{}.) Say it once more.
\end{tcolorbox}
\section[\texorpdfstring{\gu{ખાટું} (khāṭũ)}{ખાટું (khāṭũ)}]{\gu{ખાટું} (khāṭũ) — sour}

\label{lesson:GU-C132-khatu}



\begin{quote}
Before the new one: say the Gujarati for clean, then the Gujarati for dirty.

\textquotedblleft{}I think that\textquotedblright{} and \textquotedblleft{}I know that\textquotedblright{}: which small word opens the space after both? (\emph{Ke}.)
\end{quote}

\subsection*{You'll want to know: \gu{ખાટું}}
\textbf{\gu{ખાટું}} (\emph{khāṭũ}) — \textquotedblleft{}sour\textquotedblright{}.

\begin{grammarlens}[title={What you've built}]
One more word for this chapter.
\end{grammarlens}

\subsection*{Your turn}
\begin{itemize}
\raggedright
  \item \emph{Say it:} \emph{khāṭũ}
  \item \emph{Say it:} \emph{khāṭũ}, once more
  \item \emph{Say it:} say \emph{khāṭũ}
  \item \emph{Recall:} say the Gujarati for clean, then the Gujarati for dirty, then say \emph{khāṭũ} again
\end{itemize}

\begin{tcolorbox}[breakable,colback=teal!4,colframe=teal!35!black,title={Before you move on}]
What does \gu{ખાટું} mean? (\textquotedblleft{}Sour\textquotedblright{}.) Say it once more.
\end{tcolorbox}
\section[\texorpdfstring{\gu{તીખું} (tīkhũ)}{તીખું (tīkhũ)}]{\gu{તીખું} (tīkhũ) — spicy}

\label{lesson:GU-C132-tikhu}



\begin{quote}
Before the new one: say the Gujarati for dirty, then the Gujarati for sour.
\end{quote}

\subsection*{You'll want to know: \gu{તીખું}}
\textbf{\gu{તીખું}} (\emph{tīkhũ}) — \textquotedblleft{}spicy\textquotedblright{}.

\begin{grammarlens}[title={What you've built}]
One more word for this chapter.
\end{grammarlens}

\subsection*{Your turn}
\begin{itemize}
\raggedright
  \item \emph{Say it:} \emph{tīkhũ}
  \item \emph{Say it:} \emph{tīkhũ}, once more
  \item \emph{Say it:} say \emph{tīkhũ}
  \item \emph{Recall:} say the Gujarati for dirty, then the Gujarati for sour, then say \emph{tīkhũ} again
\end{itemize}

\begin{tcolorbox}[breakable,colback=teal!4,colframe=teal!35!black,title={Before you move on}]
What does \gu{તીખું} mean? (\textquotedblleft{}Spicy\textquotedblright{}.) Say it once more.
\end{tcolorbox}
\section[\texorpdfstring{\gu{નરમ} (naram)}{નરમ (naram)}]{\gu{નરમ} (naram) — soft}

\label{lesson:GU-C132-naram}



\begin{quote}
Before the new one: say the Gujarati for sour, then the Gujarati for spicy.
\end{quote}

\subsection*{You'll want to know: \gu{નરમ}}
\textbf{\gu{નરમ}} (\emph{naram}) — \textquotedblleft{}soft\textquotedblright{}.

\begin{grammarlens}[title={What you've built}]
One more word for this chapter.
\end{grammarlens}

\subsection*{Your turn}
\begin{itemize}
\raggedright
  \item \emph{Say it:} \emph{naram}
  \item \emph{Say it:} \emph{naram}, once more
  \item \emph{Say it:} say \emph{naram}
  \item \emph{Recall:} say the Gujarati for sour, then the Gujarati for spicy, then say \emph{naram} again
\end{itemize}

\begin{tcolorbox}[breakable,colback=teal!4,colframe=teal!35!black,title={Before you move on}]
What does \gu{નરમ} mean? (\textquotedblleft{}Soft\textquotedblright{}.) Say it once more.
\end{tcolorbox}
